

\section{Related Work}
\subsection{Software Engineering Agents}
% 
Inspired by the human debugging process, where developers interact with the environmental feedback (like test failures) and learn from their earlier attempts, \chatrepair~\citep{xia2023conversational,xia2024automated} proposed the first interactive bug-fixing solution based on LLMs. Since \chatrepair, a large body of research work on bug fixing and general coding tasks aims to automatically provide LLMs more context information through multi-turn conversations~\citep{chen2023teaching, kong2025contrastrepair, yuan2024evaluating}. More recently, foundation LLMs have seen substantial advances in their tool use and reasoning capabilities, and it becomes very natural to further equip feedback-driven solutions with such emergent LLM capabilities. In March 2024, Devin AI released the first AI software engineer, which can fully autonomously complete end-to-end software tasks, such as GitHub issue resolution~\citep{devin}. The initial Devin release showed an impressive resolve rate of 13.86\% on \swebench~\citep{jimenez2023swe}, a dataset with thousands of real-world GitHub issues. Since then, a large number of dedicated software agent scaffolds have been proposed, including \sweagent~\citep{yang2024swe}, \openhands~\citep{wang2024openhands}, \acr~\citep{zhang2024autocoderover}, and Trae Agent~\citep{liu2024marscode}. Such software agents typically equip the LLMs with a suite of coding tools and encourage the LLMs to autonomously decide the next actions for completing real-world software tasks. Different from the mainstream software agents, researchers have also proposed various AI software engineer solutions based on pre-defined workflows to challenge the necessity of complicated agent design, such as \agentless~\citep{xia2024agentless} and Moatless~\citep{moatless}. Moreover, recent LLMs have increasingly been post-trained on massive real-world software data to better solve software engineering issues, including SWE-RL~\cite{wei2025swerl}, DeepSWE~\cite{luodeepswe}, DeepSeek V3.1~\cite{deepseekv31}, MiniMax M1/M2~\cite{chen2025minimax}, Kimi K2~\cite{team2025kimi}, SWE-1.5~\cite{swe15}, and Code World Model (CWM)~\cite{carbonneaux2025cwm}.

Due to the huge design space for software agents, it can be extremely challenging and costly to build an optimal agent scaffold. As a result, a number of self-improving software agents have been proposed very recently, including the Self-Improving Coding Agent (SICA)~\citep{robeyns2025sica}, Darwin-Gödel Machine (DGM)~\citep{zhang2025darwin}, and Huxley-Gödel Machine (HGM)~\citep{wang2025huxley}. However, such self-improving agents require costly offline training on known benchmarks, and may not be generalize well across different LLMs, benchmarks, and issue types.
Beyond the software engineering domain, prior work~\cite{llmastool, wang2024voyager, qiu2025alitageneralistagentenabling, qian2024creatortoolcreationdisentangling, wang2024troveinducingverifiableefficient} explored using LLMs to create tools for general reasoning or embodied tasks, but they do not target real-world software engineering problems.
In contrast, in this paper, we propose \livesweagent, the first live software agent capable of performing practical self-evolution \emph{on-the-fly} during runtime when solving real-world issues. In this way, \livesweagent requires no offline training at all and can be easily generalized to different LLMs and issue domains. Moreover, it also demonstrated superior performance compared to all existing self-improving software agents. 


\subsection{Benchmarks for Software Engineering Agents}

To evaluate and demonstrate the performance of software engineering agents, a large number of datasets have been proposed. \swebench~\citep{jimenez2023swe} is one of the earliest and most widely used benchmark datasets for software agents. Besides the initial \swebench dataset, which includes thousands of real-world GitHub issues, researchers have also built curated subsets with high-quality and representative issues to support faster and more reliable evaluation, including \swebench Lite~\citep{jimenez2023swe} and \swebench Verified~\citep{sbv}. Since \swebench mostly focused on Python projects, researchers have proposed a number of \swebench-style benchmarks for projects in multiple languages for more comprehensive agent evaluation, including SWE-PolyBench~\citep{rashid2025swe}, \swebench Multilingual~\citep{sbm}, and Multi-SWE-bench~\citep{zan2025multi}. Also, \swebench relies heavily on manual effort for
collecting benchmark instances and setting up executable environments, and is hardly scalable; as a result, researchers have also leveraged software agents to streamline issue collection and environment setup to build live and scalable benchmarks, including SWE-bench-Live~\citep{zhang2025swe} and SWE-rebench~\cite{badertdinov2025swe}. More recently, Scale AI has also constructed \swebenchpro~\citep{deng2025swe}, which aims to capture more realistic, complex, enterprise-level issues than SWE-bench. For rigorous evaluation of \livesweagent, our study involved multiple widely used benchmarks, including \swebench Verified, \swebench Multilingual, \swebench Pro.

